% R15 component; only input paths and enumerated location/grammar prose are edited.
% Source: revisions/2026-10-04-r15/manuscript/model.tex; commit 1cb3cc9efe135966ec228dfaf51c6841a6dfc97c.
% Generated by code/integrate.py; edit extraction rules there.
% Reviewed source: ECTA.tex, line 75, commit f5021cefa71492babfcbfa580e0c984f59a9de26
\section{The Controlled Economy}\label{sec:model}
Let $W$ be an $m$-dimensional Brownian motion on a filtered probability space satisfying the usual conditions. The state $S=(u,X)$ contains internal and external state variables. An action $a=(\theta,c,\pi,\ldots)$ combines internal adjustment and external choices. The controlled state evolves as
\begin{equation}\label{eq:state}
 dS_s=b(s,S_s,a_s)\,ds+\sigma(s,S_s,a_s)\,dW_s,
 \qquad a_s\in\Act(s,S_s).
\end{equation}
The compact feasible correspondence $\Act$ is part of the model. Actions are progressively measurable and must generate a well-defined state process satisfying the declared state restrictions. State viability does not follow from bounded actor outputs alone.

We work initially on a finite horizon $T$. A computational domain $D$ is either invariant, equipped with a specified reflection mechanism, or stopped on exit. In the stopped case let $\tau=T\wedge\inf\{s\geq t:S_s\notin D\}$. The payoff $g(\tau,S_\tau)$ is specified on the entire parabolic stopping boundary, including the terminal face. For a fixed admissible policy $a$, utility is the policy-specific solution
\begin{equation}\label{eq:bsde}
 Y_s^a=g(\tau,S_\tau)+\int_s^\tau f(v,S_v,a_v,Y_v^a)\,dv
                    -\int_s^\tau Z_v^a\,dW_v.
\end{equation}
Adjustment costs are included in $f$. For example, time-additive utility has $f(s,S,a,y)=r_0(s,S,a)-\rho y$, where $r_0$ includes all current resource and adjustment costs. Define $V^a(t,x)=Y_t^{a,t,x}$ and $V^*(t,x)=\sup_a V^a(t,x)$. The value inside a candidate policy's recursive objective is its own utility process, not the unknown optimum inserted into every candidate objective.

\begin{assumption}[Finite-horizon comparison class]\label{ass:model}
The controlled state problem is well posed for the admissible policies used below. The coefficients are continuous, locally Lipschitz in the state uniformly over actions, and bounded on the compact stopped domain. The payoff is continuous and bounded. On an invariant utility interval containing all policy values, $f$ is continuous in its arguments and uniformly Lipschitz in utility with constant $L\geq0$. The admissible class is stable under concatenation and has the measurable-selection properties needed for dynamic programming. The resulting HJB problem and the fixed-policy problems obey comparison in the bounded continuous class with the specified boundary conditions.
\end{assumption}

The comparison requirement identifies both a function class and a boundary problem. It is not asserted merely because the primitive coefficients are continuous. On a bounded, regular Dirichlet domain, a boundary barrier is one way to verify attainment of the stopping payoff. Reflection requires the corresponding Neumann or oblique derivative condition and a well-posed reflected state process. Unbounded-state applications require growth bounds and localization. These conditions are checked separately from a neural residual.

For a smooth function $v$, write
\begin{align}
 \LL^a v&=b(t,x,a)^\top Dv+\tfrac12\tr\{\sigma(t,x,a)\sigma(t,x,a)^\top D^2v\},\label{eq:generator}\\
 \HH^a[v]&=f(t,x,a,v)+\LL^a v,\qquad
 F[v]=-v_t-\sup_{a\in\Act(t,x)}\HH^a[v].\label{eq:hjb}
\end{align}
The value solves $F[V^*]=0$ with the specified boundary data, in the viscosity sense where classical derivatives do not exist. The diffusion covariance is the positive semidefinite matrix $\sigma\sigma^\top$. No claim that the product of two arbitrary positive semidefinite matrices is symmetric positive semidefinite is needed.

A continuous martingale representation can also be used when its quadratic variation is Markov in an augmented state. If $d\langle M\rangle_s=Q(s,S_s,a_s)ds$, the covariance in \eqref{eq:generator} becomes $\sigma Q\sigma^\top$. A merely predictable process $Q_s$ is not, without further state augmentation, a deterministic coefficient of a Markov HJB equation. Quadratic variation by itself also does not imply conditionally Gaussian increments. The implementation in this paper uses the Brownian model \eqref{eq:state}.

\subsection{Infinite horizon and economic normalization}
For a stationary problem, the boundary-value equation is $\sup_a\{f(x,a,V)+\LL^aV\}=0$. We use an infinite-horizon result only after specifying admissibility, integrability, and the selected utility process. In the time-additive case, a sufficient verification requirement is that the localized It\^o identities are uniformly integrable and $\E[e^{-\rho T}|V(S_T)|]\to0$ for the policies being compared. Recursive problems require an appropriate utility comparison and transversality condition instead of simply borrowing the additive one.

Cardinal transformations also require care. Multiplying a fixed utility function by a positive constant does not change its maximizer. Adding a function of an \emph{endogenously controlled preference state} generally does. Consumption units, preference normalization, and terminal utility are therefore primitive economic choices in Section~\ref{sec:ndu}.
